% Propietario conservado: /Users/ruben/Documents/New project/output/REVISION_ARGUMENTAL_APERTURAS_CIERRES_20260915/VI/delivery/MANUSCRIPT_VI_ES_EN_COMPLETE/source_es/sections/06_kepler_sommerfeld.tex
\section{Realización central de las trayectorias y holonomía de Sommerfeld}
\label{vi:vi:sec:kepler}

La partícula definida en la sección~\ref{vi:vi:sec:particula} posee
una historia observable antes de recibir una realización orbital.
La pregunta de esta sección es cómo una representación de esa
historia conserva flujo, área y fase cuando se expresa como
movimiento central. La respuesta se organiza en cuatro operaciones:
cancelación de los flujos interiores en un corte del grafo,
normalización de su frontera en una carta métrica, evolución bajo
el operador central declarado y transporte de fase sobre sus ciclos.

El orden mantiene visible qué dato interviene en cada afirmación.
El grafo y sus registros proceden de APP--TRIT--TPK; la carta
métrica especifica cómo se mide su frontera; la realización dinámica
fija el lagrangiano que actúa en esa carta; la condición semiclásica
utiliza la constante reducida de Planck ya construida y la holonomía
de la línea transportada. Los resultados de geometría orbital se
demuestran con estos datos de partida explícitos. La mecánica
central se reconoce en la realización resultante, sin hacer de sus
constantes convencionales entradas del generador.

\subsection{Conservación de flujo y lectura de la frontera}
\label{vi:vi:subsec:gauss}

Sea $G=(V,E)$ un corte finito del grafo de rutas de la
realización considerada, con un arco en cada orientación de sus
aristas. Un flujo orientado es una función
$F:\vec E\to\mathbb R$ que satisface
$F(y,x)=-F(x,y)$. Para un vértice $x$, su divergencia es
\begin{equation}
 \operatorname{div}F(x)
 =\sum_{y:\{x,y\}\in E}F(x,y).
 \label{vi:vi:eq:divergencia-grafo}
\end{equation}
Para $S\subseteq V$, el flujo saliente se evalúa sobre los
arcos con origen en $S$ y extremo en $V\setminus S$.

\begin{theorem}[Balance de flujo en un corte finito]
\label{vi:vi:thm:gauss}
Para toda parte $S\subseteq V$,
\begin{equation}
 \sum_{x\in S}\operatorname{div}F(x)
 =\sum_{\substack{x\in S,\ y\notin S\\\{x,y\}\in E}}F(x,y).
 \label{vi:vi:eq:gauss}
\end{equation}
\end{theorem}
\begin{proof}
Una arista con ambos extremos en $S$ contribuye dos veces a la
suma de divergencias, mediante $F(x,y)+F(y,x)=0$. Las aristas
con ambos extremos fuera de $S$ no contribuyen. Quedan
exactamente las aristas que atraviesan la frontera, con la
orientación indicada.
\end{proof}

El teorema transmite a la frontera el balance producido en el
interior. Para una representación cúbica con vecinos inmediatos,
considérese
\[
 C_R=\{x\in\mathbb Z^3:\|x\|_\infty\leq R\},
 \qquad R\in\mathbb Z_{\geq0}.
\]
Cada una de las seis caras contiene $(2R+1)^2$ enlaces
salientes en su dirección normal. Los enlaces de direcciones
distintas se cuentan por separado, incluso cuando parten de
un mismo vértice de arista o esquina. Por tanto,
\begin{equation}
 |\partial C_R|=6(2R+1)^2.
 \label{vi:vi:eq:frontera-cubica}
\end{equation}
Este censo corresponde a enlaces salientes, no al número de
vértices de la frontera.

Si $Q_F$ es el flujo total de la fuente encerrada, su media
por enlace es $Q_F/[6(2R+1)^2]$. Bajo la condición de isotropía
media del lector, esa media es la intensidad que se asigna al
cascarón. Definamos su radio efectivo, en unidades del paso de
la carta, por
\begin{equation}
 4\pi r_{\mathrm{eff}}^2=6(2R+1)^2.
 \label{vi:vi:eq:radio-efectivo}
\end{equation}
Se obtiene entonces
\begin{equation}
 E_R=\frac{Q_F}{4\pi r_{\mathrm{eff}}^2}.
 \label{vi:vi:eq:inverso-area}
\end{equation}
La coordenada $\pi$ de esta normalización es la salida regional
HMT ya publicada. La identidad procede del censo y de la
definición del radio; el paso de una media a una realización
central utiliza, además, la isotropía indicada. La carta
dimensional decide después si $Q_F$ representa flujo eléctrico,
gravitatorio u otro observable y cómo convierte el paso
reticular en longitud. El balance finito conserva ese
distinto estatuto de cada operación.

\subsection{Dinámica central en la hoja euclídea}
\label{vi:vi:subsec:dinamica-central}

Sea $\mu>0$ la masa reducida de la realización material y sea
$\mathcal K>0$ su acoplamiento central, con dimensión
producto de energía y longitud. Estas coordenadas se reciben después
de las construcciones internas de masa, acción y acoplamiento;
su origen y su recta dimensional se conservan al efectuar la
realización. El teorema siguiente tiene como dominio todas
las realizaciones que, en su hoja euclídea, adopten
\begin{equation}
 L(\mathbf r,\dot{\mathbf r})
 =\frac\mu2\|\dot{\mathbf r}\|^2+\frac{\mathcal K}{r},
 \qquad
 H(\mathbf r,\mathbf p)
 =\frac{\|\mathbf p\|^2}{2\mu}-\frac{\mathcal K}{r},
 \quad r=\|\mathbf r\|>0,
 \quad\mathbf p=\mu\dot{\mathbf r}.
 \label{vi:vi:eq:kepler-lh}
\end{equation}
La exclusión de $r=0$ delimita el dominio regular de estas
coordenadas. La trayectoria se considera en un intervalo de
tiempo sin colisión.

\begin{proposition}[Plano orbital y ecuación de la cónica]
\label{vi:vi:prop:conica}
Una solución no radial de la ecuación de Euler--Lagrange de
\eqref{vi:vi:eq:kepler-lh} conserva
$\boldsymbol\ell=\mathbf r\times\dot{\mathbf r}$, con
$\ell=\|\boldsymbol\ell\|>0$. En el plano ortogonal a
$\boldsymbol\ell$ puede escribirse
\begin{equation}
 r(\theta)=\frac{p}{1+\epsilon\cos(\theta-\theta_0)},
 \qquad p=\frac{\mu\ell^2}{\mathcal K},
 \qquad \epsilon\geq0.
 \label{vi:vi:eq:conica}
\end{equation}
La energía negativa corresponde al régimen elíptico
$0\leq\epsilon<1$.
\end{proposition}
\begin{proof}
La derivada del potencial da
\begin{equation}
 \mu\ddot{\mathbf r}=-\frac{\mathcal K}{r^3}\mathbf r.
 \label{vi:vi:eq:fuerza-central}
\end{equation}
De aquí,
$\dot{\boldsymbol\ell}=\mathbf r\times\ddot{\mathbf r}=0$.
La posición y la velocidad permanecen ortogonales al vector
constante no nulo $\boldsymbol\ell$, lo que fija el plano.
En sus coordenadas polares $r^2\dot\theta=\ell$. Al introducir
$u=1/r$ en la ecuación radial se obtiene la ecuación de Binet
\[
 u''(\theta)+u(\theta)=\frac{\mathcal K}{\mu\ell^2}.
\]
Su solución general es
$u=\mathcal K[1+\epsilon\cos(\theta-\theta_0)]/(\mu\ell^2)$,
y su inversión prueba \eqref{vi:vi:eq:conica} en el dominio
donde el denominador es positivo. Sustituir la solución en
la energía da
$H=\mathcal K(\epsilon^2-1)/(2p)$, que es negativa
exactamente cuando $\epsilon<1$.
\end{proof}

La letra $\epsilon$ designa aquí la excentricidad orbital y
no el número de Euler. El vector $\boldsymbol\ell$ es un
momento angular específico, no la constante de Planck.
Estas elecciones de notación permiten comparar, sin mezclarlos,
los observables geométricos y las constantes internas que
intervienen en su realización.

\begin{proposition}[Conservación areal y ley del periodo]
\label{vi:vi:prop:ley-areal}
La velocidad areolar de la órbita es constante:
\begin{equation}
 \frac{d\mathcal A_{\mathrm{orb}}}{dt}=\frac\ell2.
 \label{vi:vi:eq:ley-areal}
\end{equation}
Para una elipse de semieje mayor $a$, su periodo satisface
\begin{equation}
 T^2=\frac{4\pi^2\mu}{\mathcal K}a^3.
 \label{vi:vi:eq:tercera-ley}
\end{equation}
\end{proposition}
\begin{proof}
El sector barrido en tiempo $dt$ tiene área
$r^2d\theta/2=\ell\,dt/2$. Si $b=a\sqrt{1-\epsilon^2}$
es el semieje menor, el área total es $\pi ab$ y
$T=2\pi ab/\ell$. La relación
$p=a(1-\epsilon^2)=\mu\ell^2/\mathcal K$ implica
$\ell^2=\mathcal K a(1-\epsilon^2)/\mu$.
Sustituir en $T^2=4\pi^2a^2b^2/\ell^2$ da la expresión
\eqref{vi:vi:eq:tercera-ley}.
\end{proof}

La conservación areal está determinada por la simetría central
del movimiento. La elipse de esta proposición vive en el
espacio de configuración. La elipse de acción construida
desde el par angular vive, en cambio, en su carta de acción
y transporta las normalizaciones correspondientes. El enlace
entre ambas exige el mapa de realización que identifica sus
variables; compartir una forma elíptica no sustituye ese mapa.
La lectura orbital conserva así las dependencias de masa,
acoplamiento y geometría que la producen.

\subsection{Cubierta orbital y línea real de signo}
\label{vi:vi:subsec:cubierta-orbital}

En el plano complejo sin el origen, la aplicación
\begin{equation}
 q=z^2:\mathbb C\setminus\{0\}\longrightarrow
                 \mathbb C\setminus\{0\}
 \label{vi:vi:eq:levi-civita}
\end{equation}
es una cubierta doble. Si $q(t)=r(t)e^{i\theta(t)}$ y se
elige inicialmente una raíz, su levantamiento es
$z(t)=\sqrt{r(t)}e^{i\theta(t)/2}$. Una vuelta con incremento
$2\pi$ de $\theta$ termina en $-z(0)$, y dos vueltas
restituyen el punto inicial. Este es el transporte de los dos
levantamientos sobre el plano perforado.

Una banda orbital requiere otro objeto: una línea real
asociada a la órbita y a su carácter de signo. Sea
$\gamma:S^1\to\mathbb R^2$ una parametrización regular
inyectiva de una órbita elíptica y sea $\delta>0$. Con
holonomía $-1$ en una vuelta, su banda de intervalos se
representa por
\begin{equation}
 M_\gamma=([0,2\pi]\times[-\delta,\delta])/
          \bigl((0,u)\sim(2\pi,-u)\bigr).
 \label{vi:vi:eq:mobius-orbital}
\end{equation}
La identificación de la base con la imagen de $\gamma$
sitúa la sección nula sobre la órbita.

\begin{proposition}[Transporte orientado sobre el ciclo orbital]
\label{vi:vi:prop:holonomia-orbital}
El transporte de la línea de \eqref{vi:vi:eq:mobius-orbital}
satisface
$\mathcal H(\gamma)=-I$ y $\mathcal H(\gamma^2)=I$.
\end{proposition}
\begin{proof}
La fibra terminal se identifica con la inicial por
$u\mapsto-u$; componer ese mapa dos veces da la identidad.
Es la construcción del
teorema~\ref{vi:vi:thm:retorno-mobius}, aplicada al lazo orbital.
\end{proof}

La cubierta \eqref{vi:vi:eq:levi-civita} tiene fibras de dos
puntos; \eqref{vi:vi:eq:mobius-orbital} tiene fibras de
intervalos. La monodromía de orden dos aparece en ambas,
pero conserva esos tipos distintos. La regularización
dinámica de Levi--Civita utiliza, además, la transformación
de momentos y la reparametrización temporal. En la presente
construcción se utiliza únicamente su cubierta y su ley de
levantamiento, con el alcance que acaba de demostrarse.

\subsection{Condición de fase y desplazamiento semientero}
\label{vi:vi:subsec:sommerfeld}

Sea $\gamma$ un ciclo de la realización hamiltoniana y sea
\begin{equation}
 J_\gamma=\oint_\gamma\mathbf p\cdot d\mathbf q
 \label{vi:vi:eq:accion-ciclo}
\end{equation}
su acción. La constante reducida de Planck
$\hbar_{\mathrm{int}}>0$ se recibe de la construcción HMT
de acción, en la misma recta dimensional que $J_\gamma$.
El cociente $J_\gamma/\hbar_{\mathrm{int}}$ es, por tanto,
adimensional. Supóngase que el transporte semiclásico en la
línea complejificada aporta la fase
$\exp(iJ_\gamma/\hbar_{\mathrm{int}})$ y que la línea real
de orientación aporta el carácter
$\varepsilon_\gamma\in\{+1,-1\}$. La condición de cierre
de la sección transportada es
\begin{equation}
 \varepsilon_\gamma
       \exp\!\left(\frac{iJ_\gamma}{\hbar_{\mathrm{int}}}\right)=1.
 \label{vi:vi:eq:cierre-sommerfeld}
\end{equation}
El producto conserva las dos contribuciones: acción del
ciclo y signo de la línea. La segunda no se absorbe en una
modificación de la constante de Planck.

\begin{theorem}[Sectores entero y semientero de la acción]
\label{vi:vi:thm:sommerfeld}
Bajo el transporte \eqref{vi:vi:eq:cierre-sommerfeld},
\begin{equation}
 J_\gamma=
 \begin{cases}
  2\pi\hbar_{\mathrm{int}}n,
          &\varepsilon_\gamma=+1,\\[1mm]
  2\pi\hbar_{\mathrm{int}}(n+\tfrac12),
          &\varepsilon_\gamma=-1,
 \end{cases}
 \qquad n\in\mathbb Z.
 \label{vi:vi:eq:accion-sommerfeld}
\end{equation}
\end{theorem}
\begin{proof}
Para signo positivo, la exponencial debe valer $1$, cuyos
argumentos reales son $2\pi n$. Para signo negativo, debe
valer $-1$, cuyos argumentos son $(2n+1)\pi$. Multiplicar
por $\hbar_{\mathrm{int}}$ produce las dos familias.
\end{proof}

La fórmula determina la restricción de cierre sobre la
acción de un ciclo dado. Para obtener un espectro atómico
se utiliza, además, el Hamiltoniano específico, los ciclos
independientes y la prescripción de cuantización. Si esta
prescripción incluye índices de Maslov, sus fases se
componen con \eqref{vi:vi:eq:cierre-sommerfeld} conforme al
transporte escogido, evitando contar dos veces una misma
contribución topológica.

Invertir la orientación del ciclo cambia $J_\gamma$ por
$-J_\gamma$, mientras el carácter de signo cumple
$\varepsilon_{\gamma^{-1}}=\varepsilon_\gamma^{-1}
=\varepsilon_\gamma$. La fase total se transforma en su
inversa y la condición de cierre se conserva. Esta es una
ley precisa de bidireccionalidad orbital. Su identificación
con las antisecciones materiales requiere el levantamiento
de la sección~\ref{vi:vi:subsec:antisecciones}; la igualdad
de dos signos no establece por sí sola esa identificación.

La realización reúne, de este modo, tres conservaciones
consecutivas: el flujo pasa del interior a la frontera;
la dinámica central conserva el momento angular y el área
barrida por unidad de tiempo; el transporte de la sección
conserva su condición de cierre mediante acción y
holonomía. Cada resultado conserva la ecuación, el dominio
y los datos que lo sostienen. La historia material del
TPK se representa en esa dinámica con su memoria, en lugar
de quedar sustituida por la figura orbital que publica.
